%! TEX root = ../main.tex
\documentclass[../main.tex]{subfiles}
\begin{document}

\subsection {Công nghệ phía máy chủ}

Về phía máy chủ, nhóm đã thống nhất sử dụng ngôn ngữ lập trình Python để phát triển ứng dụng. Đây là một ngôn ngữ lập trình rất phổ biến hiện nay trong phát triển ứng dụng web.
Nó có cấu trúc mã nguồn rất đơn giản và dễ đọc, giúp người lập trình có thể làm quen nhanh hơn và tập trung vào phát triển hệ thống thay vì phải học cú pháp phức tạp. Đồng thời
điều đó còn giúp cả nhóm phát triển đều hiểu được chức năng dù không tham gia lập trình trực tiếp, nhằm tăng hiệu quả phối hợp. Ngoài ra, Python còn có một cộng đồng đông đảo các nhà phát triển,
cung cấp lượng lớn các thư viện hỗ trợ nhiều nhu cầu khác nhau cho người dùng. Những điểm trên đã giúp cho quá trình phát triển và triển khai ứng dụng trở nên nhanh chóng và dễ dàng hơn.

Nhóm cũng đã lựa chọn framework FastAPI để xây dựng API cho ứng dụng. Đây là một framework hiện đại sử dụng ngôn ngữ Python. Đúng như tên gọi, thư viện này hỗ trợ việc tạo nên các API cực kì nhanh chóng với
cú pháp đơn giản. Đây cũng là một trong những framework nhanh nhất của Python hiện nay, với tốc độ ngang bằng với các ngôn ngữ khác như NodeJS hay Golang chuyên dụng cho phát triển ứng dụng web. Vì thế nó phù hợp
để xây dựng lên hệ thống diễn đàn cần xử lí lượng lớn các yêu cầu. Các API được tạo ra bởi FastAPI cũng tuân thủ chuẩn OpenAPI, giúp việc tích hợp vào ứng dụng khách trở nên dễ dàng và có hệ thống.
FastAPI còn giúp tạo ra tài liệu API tự động có thể tương tác và kiểm thử trực tiếp trên trình duyệt một cách trực quan. Nó cho phép các đội nhóm phát triển khác sử dụng thử, và cung cấp tài liệu hướng dẫn cài đặt và sử dụng API của hệ thống.
Ngoài ra, nó còn cung cấp các tính năng cơ bản mà ứng dụng web cần có như xác thực và bảo mật, kiểm tra dữ liệu đầu vào, và nhiều tính năng khác.

Đối với cơ sở dữ liệu, nhóm đã lựa chọn PostgreSQL. Đây là một cơ sở dữ liệu quan hệ mã nguồn mở phổ biến hiện nay. Chức năng của nó là lưu trữ, quản lí và tổ chức dữ liệu một cách có hệ thống dưới dạng các bảng.
Mỗi bảng lại bao gồm các trường, hay các cột dữ liệu thể hiện các thuộc tính của đối tượng được lưu, và các bản ghi, hay các hàng dữ liệu là một dữ liệu cụ thể cho một đối tượng. Các bảng có thể liên kết với nhau thông qua các trường liên kết dữ liệu, gọi là khóa ngoại.
Điều này giúp thể hiện các quan hệ của các đối tượng trong hệ thống. PostgreSQL phù hợp với điều kiện sử dụng phức tạp của diễn đàn, với các đối tượng quan hệ đa dạng. Và vì nó sử dụng ngôn ngữ SQL thuần túy nên rất dễ làm quen.

Để kết nối và tương tác với cơ sở dữ liệu, hệ thống sử dụng thư viện SQLAlchemy. Đây là một bộ công cụ ORM giúp ánh xạ các thuộc tính đối tượng trong hệ thống tương ứng với các trường trong cơ sở dữ liệu. Từ đó việc thao tác đọc ghi dữ liệu trở nên nhẹ nhàng hơn thông qua việc thao tác trên chính các đối tượng.
Ngoài ra SQLAlchemy còn hỗ trợ tạo các truy vấn phức tạp, hỗ trợ truy cập đến các đối tượng quan hệ như là các thuộc tính của đối tượng cha thay vì phải viết biểu thức SQL.


Ngoài ra, để tăng tốc độ truy xuất dữ liệu, hệ thống sử dụng Redis làm bộ nhớ đệm. Redis cũng là một cơ sở dữ liệu lưu trữ dưới dạng các cặp khóa - giá trị. Redis được tổ chức để hoạt động trên bộ nhớ RAM nên có tốc độ truy xuất cực kì ấn tượng, rất phù hợp với các thao tác đọc dữ liệu lặp lại nhiều lần như trong hệ thống diễn đàn.
Redis hiện nay được sử dụng rộng rãi trong các hệ thống lớn, nên cũng có cộng đồng hỗ trợ đông đảo và tài liệu hướng dẫn chi tiết. Ngoài ra Redis còn hỗ trợ cấu hình Pub/Sub, rất thích hợp cho hoạt động thông báo được sử dụng trong diễn đàn.

Bên cạnh đó, hệ thống cũng sử dụng dịch vụ SMTP của Gmail để thực hiện chức năng gửi thư điện tử đến người dùng. Gmail cung cấp dịch vụ SMTP miễn phí và rất dễ sử dụng. Cùng với sự tin cậy và tốc độ gửi thư nhanh thì đây là một lựa chọn hợp lí cho ứng dụng.

Cuối cùng, để hệ thống được vận hành ổn định, nhóm đã lựa chọn sử dụng dịch vụ điện toán đám mây Azure của Microsoft. Đây là một trong những nhà cung cấp dịch vụ đám mây lớn nhất hiện nay, với hệ sinh thái dịch vụ đa dạng và phong phú. Các dịch vụ mà hệ thống có thể khai thác bao gồm: triển khai ứng dụng web với Azure Web app, cơ sở dữ liệu với Azure Database for PostgreSQL flexible server, bộ nhớ đệm với Azure managed Redis, và kho lưu trữ dữ liệu đa phương tiện Azure Blob Storage.
Điều này giúp việc quản lí vận hành hệ thống máy chủ đơn giản hơn, đảm bảo chất lượng hoạt động tốt với thời gian hoạt động cao, dễ dàng mở rộng, và tiết kiệm chi phí phần cứng cũng như bảo trì.

\end{document}